% !PRINTED 228
\weberpage{228}
Die Anzahl der Di\"edergruppen erhalten wir daraus, dass die Di\"edergruppe durch die in ihr enthaltene cyklische Gruppe vollständig bestimmt ist (\S.\,55). Bei der Vierergruppe ist aber noch zu beachten, dass man dieselbe Gruppe erhält, wenn man von $\psi$, von $\chi$ oder von $\psi\chi$ ausgeht, und dass also die direct erhaltene Zahl durch 3 zu dividiren ist. Die Anzahl der $D_2$ ist also 5, die der $D_3$ ist 10 und die der $D_5$ ist 6.

Von besonderer Wichtigkeit sind aber die in der Ikosa\"edergruppe enthaltenen Tetra\"edergruppen. Wir wollen eine von ihnen, die wir mit $Q$ bezeichnen, bestimmen.

Die Tetra\"edergruppe muss (nach \S.\,56) eine Vierergruppe $D_2$ enthalten. Wir gehen von einer solchen Gruppe $D_2$ aus und wählen dazu
\[
1, \psi, \chi, \psi\chi.
\]

Es muss nun weiter in $Q$ eine Substitution dritter Ordnung vorkommen. Diese können wir in einer der beiden Formen $\Theta^r\chi\Theta^s$ oder $\Theta^r\chi\psi\Theta^s$ annehmen. Da beide Annahmen zu demselben Resultate führen, wählen wir als Substitution dritter Ordnung
\[
\varphi = \Theta^r\chi\Theta^s, \quad r+s\equiv\pm2 \ (\mathrm{mod}\ 5) \quad [\S.\,58, (21)].
\]

Von den Zahlen $r, s$ kann aber keine $\equiv0$ sein, weil sonst entweder $\chi\varphi$ oder $\varphi\chi$ eine Potenz von $\Theta$, also vom $5^{\mathrm{ten}}$ Grade wäre, während doch in $Q$ kein Element $5^{\mathrm{ten}}$ Grades vorkommen kann.

Wir bilden ferner nach \S.\,58, (23), (25):
\[
\begin{aligned}
\chi\varphi &= \Theta^r\chi\psi\Theta^{r+s}, \quad r\equiv\pm1 \\
\chi\varphi &= \Theta^{-r}\chi\Theta^{-r+s}, \quad r\equiv\pm2,
\end{aligned}
\]
wodurch aus dem oben angeführten Grunde im ersten Falle $r\equiv-s$, im zweiten $r\equiv s$ ausgeschlossen ist, und im ersten Falle $2r+s\equiv\pm1$, im zweiten $-2r+s\equiv\pm2$ gefordert wird. Danach bleiben die vier möglichen Fälle
\[
r\equiv\pm1, \quad s\equiv\pm2, \quad r\equiv\pm2, \quad s\equiv\pm1 \ (\mathrm{mod}\ 5)
\]
übrig.

Wenn von den so bestimmten vier Substitutionen dritter Ordnung eine in $Q$ vorkommt, so kommen auch die drei anderen vor. Denn setzen wir
\weq{1}{\varphi = \Theta\chi\Theta^2,}
so ergiebt sich nach \S.\,58, (23), (25):
\weq{2}{\varphi^{-1} = \Theta^{-2}\chi\Theta^{-1}, \quad \varphi\chi = \Theta^{-1}\chi\Theta^{-2}, \quad \chi\varphi^{-1} = \Theta^2\chi\Theta.}
